% demo1-fixed：demo1.tex 的 NTex 自举版（当前引擎为 iniTeX 初始态、无 plain 格式预载，
% 需显式装配字体族与 plain 宏；对照 demo1.tex 的依赖缺口）。
\hsize=6.5in
\vsize=9.0in

% ---- 字体装配（plain.tex 等价手写）----
\font\tenrm=cmr10
\font\titlefont=cmr12 scaled 1728
\font\boldfont=cmbx10 scaled 1200
\font\tenti=cmti10
\font\tentt=cmtt10
\font\tenmi=cmmi10
\font\tensy=cmsy10
\font\tenex=cmex10
\font\sevenrm=cmr7
\font\sevenmi=cmmi7
\font\sevensy=cmsy7
\font\fiverm=cmr5
\font\fivemi=cmmi5
\font\fivesy=cmsy5
\textfont0=\tenrm
\scriptfont0=\sevenrm
\scriptscriptfont0=\fiverm
\textfont1=\tenmi
\scriptfont1=\sevenmi
\scriptscriptfont1=\fivemi
\textfont2=\tensy
\scriptfont2=\sevensy
\scriptscriptfont2=\fivesy
\textfont3=\tenex
\tenrm

% ---- plain 关键垂直/行距参数（无格式预载时的等价手写）----
\topskip=10pt
\parskip=0pt
\parindent=0pt
\baselineskip=12pt
\lineskip=1pt
\lineskiplimit=0pt
\abovedisplayskip=12pt plus3pt minus9pt
\belowdisplayskip=12pt plus3pt minus9pt
\abovedisplayshortskip=0pt plus3pt
\belowdisplayshortskip=7pt plus3pt minus4pt

% ---- 数学命令（plain.tex mathchardef 摘录）----
\mathchardef\sum="1350
\mathchardef\infty="1231
\mathchardef\pi="0170
\mathchardef\bullet="220F

% ---- plain 宏最小集 ----
\def\centerline#1{\hbox to\hsize{\hss#1\hss}}
\def\medskip{\vskip6pt plus2pt minus2pt}
\def\it{\tenti}
\def\tt{\tentt}
\def\item#1{\par\hangindent=1.5em\hangafter=1\noindent\hbox to1.5em{#1\hss}\ignorespaces}
\def\LaTeX{LaTeX}
\def\bye{\par\vfill\eject\end}

\centerline{\titlefont Plain TeX Demonstration}
\vskip 0.5in

This is a simple document typeset using standard {\boldfont Plain TeX} rather
than LaTeX. In Plain TeX, paragraph breaks are created by leaving a blank line
in the source file, just like in LaTeX.

\vskip 0.25in

{\boldfont 1. Mathematical Formulas}
\medskip
You can write inline math like $E = mc^2$, or displayed math equations using
double dollar signs:
$$ \sum_{n=1}^{\infty} {1 \over n^2} = {\pi^2 \over 6} $$
Notice that instead of \LaTeX's \vbox{\tt\string\frac\char`\{1\char`\}\char`\{n\char`\^2\char`\}}, Plain TeX uses the
native {\tt\string\over} primitive.

\vskip 0.25in

{\boldfont 2. Bulleted Lists}
\medskip
Plain TeX uses the {\tt\string\item} macro for lists. It does not require wrapping
them inside an environment:

\item{$\bullet$} This is the first item in a native TeX list.
\item{$\bullet$} This is the second item.
\item{$\bullet$} Notice there is no closing tag needed for this list.

\vskip 0.5in
\centerline{\it End of Document.}

\bye
